% Please do not change %
\documentclass[11pt]{article} % font size
\usepackage[margin=1in]{geometry} % margins
\usepackage{amsmath,amsthm,amssymb} % for the  common "math symbols"
\usepackage[shortlabels]{enumitem} % for enumeration
\usepackage{booktabs} % if you need tables
\usepackage{listings}
\usepackage{color}

\lstset{
  frame=none,
  xleftmargin=2pt,
  stepnumber=1,
  numbers=left,
  numbersep=5pt,
  numberstyle=\ttfamily\tiny\color[gray]{0.3},
  belowcaptionskip=\bigskipamount,
  captionpos=b,
  escapeinside={*'}{'*},
  language=haskell,
  tabsize=2,
  emphstyle={\bf},
  commentstyle=\it,
  stringstyle=\mdseries\rmfamily,
  showspaces=false,
  keywordstyle=\bfseries\rmfamily,
  columns=flexible,
  basicstyle=\small\sffamily,
  showstringspaces=false,
  morecomment=[l]\%,
}
% Feel free to include additional packages (if needed), but make sure it does not override the margin/font requirements.
%
\setlength{\parindent}{0pt} % no indent for new paragraphs.
\setlength{\parskip}{5pt} % distance between paragraphs, which will be used when you have a blank line in your LaTeX code.


%%%%%%%%%%%%%%%%%%%%%%%%%%%%%%%%%%%%%%%%%%%%%%%%%%%%%%%%%%%%%%%%%%%%%%%%%%%%%%
\title{Principles of Programming Languages: \\Problem Sheet 3 - Semantic variations}
%
\author{Wira Azmoon Ahmad}
\date{17 Nov 2023}
%
\begin{document}
%
\maketitle

%You may use the usual \LaTeX\ environments to structure your answers:
%\begin{align}
%        \label{basegnn}
%        \mathbf{h}_u^{(t+1)} = \sigma \bigg( \mathbf{W}_\text{self}^{(t)} \mathbf{h}_u^{(t)}  + \mathbf{W}_\text{neigh}^{(t)} \sum_{v \in N(u)} \mathbf{h}_v^{(t)} + \mathbf{b}^{(t)} \bigg). 
%\end{align}


%\begin{table}[h]
%\caption{A simple table.}
%\centering
%\begin{tabular}[t]{lrrrr}
%\toprule
%Graphs & Can & Be  & Fun & Too \\ 
%\midrule 
%$G_1$ & $1$  & $2$ & $3$ & $4$ \\ 
%$G_2$ & $-1$  & $-2$ & $-3$ & $-4$ \\ 
%$G_3$ & $0.1$  & $0.2$ & $0.3$ & $0.4$ \\ 
%\bottomrule
%\end{tabular}
%\label{tab}
%\end{table}

\section*{Question 1}

% To answer each part of the question, please simply use a paragraph command 
% \paragraph{(a)}
\lstinputlisting[language=Haskell, lastline=30]{../PS3.hs}

\clearpage
\section*{Question 2}

\paragraph{(a)}

\lstinputlisting[language=Haskell, linerange={11-17}]{../FunMonadExnVal.hs}

\paragraph{(b)}

\lstinputlisting[language=Haskell, linerange={19-25}]{../FunMonadExnVal.hs}

\paragraph{(c)}

\lstinputlisting[language=Haskell, linerange={33-35}]{../PS3.hs}

\paragraph{(d)}
The expression gives 2. We can force the \texttt{eval e2 env} before \texttt{orelse}.

\lstinputlisting[language=Haskell, linerange={37-38}]{../PS3.hs}

\section*{Question 3}
\lstinputlisting[language=Haskell, linerange={41-73}]{../PS3.hs}

\section*{Question 4}
\lstinputlisting[language=Haskell, linerange={75-103}]{../PS3.hs}

\section*{Question 5}
\paragraph{(a)}
$M_1$ $\alpha$ maps the memory to a type $Maybe$, applied on a tuple of $(\alpha,Mem)$
while $M_2$ $\alpha$ maps the memory to a tuple, within which there's a $Maybe$ and $Mem$ type.

\paragraph{(b)}
This means that the first type keeps memory only if the $Maybe$ type is $Just$, not if it is $Nothing$,
while the second type always keeps memory whether or not the $Maybe$ type is $Just$ or $Nothing$.
The first type is more efficient since it does not worry about memory as much.

\paragraph{(c)}
The example below has 1 for the first interpreter, and 2 for the second.
\lstinputlisting[lastline=3]{../PS3.fun}

\section*{Question 6}
\lstinputlisting[linerange={65-67,88-101,105-112}]{../FunC.hs}

\clearpage
\section*{Question 7}
\paragraph{(a)}

\lstinputlisting[language=Haskell]{../swapmem.fun}

% \paragraph{(b)}
\lstinputlisting[language=Haskell,linerange={102-104}]{../FunC.hs}

% \paragraph{(c)}
\lstinputlisting[language=Haskell]{../swapc.fun}

% \paragraph{(d)}
\lstinputlisting[linerange={122-126}]{../PS3.hs}

\paragraph{(e)}
No, since $x$ may not be an address, so dereferencing it does not make sense.
\lstinputlisting[language=Haskell,linerange={128-133}]{../PS3.hs}

\clearpage
\section*{Question 8}
\lstinputlisting[language=Haskell,linerange={136-138}]{../PS3.hs}

\section*{Question 9}
\lstinputlisting[language=Haskell,firstline=140]{../PS3.hs}


\end{document}